\documentclass{article}
\usepackage[utf8]{inputenc}

\usepackage[a4paper, total={6in, 8in}]{geometry}

\usepackage{amsfonts} 
\def\sone{{\textbf{1}}}

\usepackage{graphicx}
\usepackage{amsmath}
\usepackage{amssymb}
\usepackage{bbm}
\usepackage{xcolor}
\usepackage{hyperref}
\hypersetup{
    colorlinks=true,
    linkcolor=blue,
    filecolor=magenta,      
    urlcolor=cyan,
    pdftitle={Overleaf Example},
    pdfpagemode=FullScreen,
}

\title{\textbf{Assignment 1}}
\author{\textbf{\color{red} {Deadline: Tuesday, $5^{th}$ of May 2026}}}
\date{\textbf{\color{red} Upload your solutions at: \url{https://tinyurl.com/AML-2026-ASSIGNMENT1}
}}

\begin{document}

\maketitle

1. \textbf{(0.3 points)} What is the maximum value of the natural even number $n$, $n=2m$, such that there exists a hypothesis class $\mathcal{H}$ with $n$ elements that shatters a set C of $m = \frac{n}{2}$ points? Give an example of such an $\mathcal{H}$ and C. Justify your answer.




\vspace{5mm}

2. \textbf{(1.2 points)} Consider $\mathcal{H} = \mathcal{H}_1 \cup \mathcal{H}_2 \cup \mathcal{H}_3$, where:

\begin{itemize}
    \setlength{\itemindent}{-6mm}
    \item[] $\mathcal{H}_1 = \{h_{a} : \mathbb{R} \rightarrow \{0,1\} \ | \ h_{a}(x) = \sone_{[x \geq a]}(x) = \sone_{[a, +\infty)}(x), a \in \mathbb{R}\}$,
    \item[] $\mathcal{H}_2 = \{h_{a} : \mathbb{R} \rightarrow \{0,1\} \ | \ h_{a}(x) = \sone_{[x < a]}(x) = \sone_{(-\infty, a)}(x), a \in \mathbb{R}\}$,
    \item[] $\mathcal{H}_3 = \{h_{a, b} : \mathbb{R} \rightarrow \{0,1\} \ | \ h_{a, b}(x) = \sone_{[a \leq x \leq b]}(x) = \sone_{[a, b]}(x), a, b \in \mathbb{R}\}$.
\end{itemize}    

Consider the realizability assumption.

\begin{itemize}
    \item[a)] Compute VCdim($\mathcal{H}$).
    \item[b)] Give an algorithm A and determine an upper bound on the sample complexity $m_{\mathcal{H}}(\epsilon, \delta)$ such that the definition of PAC-learnability is satisfied.
    
\end{itemize}



3. \textbf{(1.2 points)} Usually, there is a correlation between the number of parameters of a hypothesis class $\mathcal{H}$ and its VC-dimension. Consider the following hypothesis classes:

\begin{itemize}
    \setlength{\itemindent}{-6mm}
    \item[] $\mathcal{H}_1 = \{h_{a} : \mathbb{R} \rightarrow \{0,1\} \ | \ h_{a}(x) = \sone_{[a, a+1]}(x), a \in \mathbb{R}\}$,
    \item[] $\mathcal{H}_2 = \{h_{b} : \mathbb{R} \rightarrow \{0,1\} \ | \ h_{b}(x) = \sone_{[b,b+1] \cup [b+2,b+4]}(x), b \in \mathbb{R}\}$,
    \item[] $\mathcal{H}_3 = \{h_{c} : \mathbb{R} \rightarrow \{0,1\} \ | \ h_{c}(x) = \sone_{[c,c+1] \cup [c+2,c+4] \cup [c+6,c+9]}(x), c \in \mathbb{R}\}$.
\end{itemize}    

Consider the realizability assumption.  Compute VCdim($\mathcal{H}_1$), VCdim($\mathcal{H}_2$), VCdim($\mathcal{H}_3$).

\vspace{5mm}


\textbf{Ex-officio: 0.3 points}

\end{document}